\chapter{因子分析}
因子分析与主成分分析在计算上相似，均需求解特征值，但二者的方向相反。主成分分析由变量构造综合指标，$c=a\trans x$，主成分为变量的加权和，数据给定即可计算。因子分析由因子解释变量，$x=Lf+u$，假定存在不可观测的公因子，变量之间的相关源于公因子的共同作用。

因此，因子分析以一种解释为出发点：若干指标之间相关，是因为它们有共同的来源。这与流行病学中两个结局因共同危险因素而相关的思路一致。在心理测量中，言语能力、数学能力等不可直接观测，只能通过各科成绩加以推断，因子分析即用于此类推断。

本章的重点为：模型$\Sigma=LL\trans+\Psi$的由来；负荷、共性方差与因子贡献的含义；因子旋转不改变模型拟合而改变解释的原因。
\section{因子分析}
因子分析（factor analysis）探讨多个可观测变量间的相关性，解释相关阵或协方差阵为什么会分块。有研究者认为，$p$个表型变量受少数$k$个（$k<p$）不可观测隐变量（因子）支配。因子分析的任务是探讨因子如何支配变量，还原它们之间的协方差结构。
\Note{主成分分析是数学变换，目的为降维与信息浓缩。因子分析是统计模型，目的是寻找潜在变量（公因子）并解释它们的意义，且允许误差。}
\subsection{发展历程}
1904年，Charles Spearman认为学生古典语、法语和英语三门成绩受公因子“语言能力”支配，用一个不可观测变量解释可观测变量间的相关性（协方差结构）。
\[
R=\begin{array}{c@{\quad}c}
\begin{matrix}\text{古典语}\\\text{法语}\\\text{英语}\end{matrix}&
\begin{bmatrix}1.00&0.83&0.78\\0.83&1.00&0.67\\0.78&0.67&1.00\end{bmatrix}
\end{array}.
\]
\begin{sikao}[单因子模型的相关结构]
若三门成绩仅受一个公因子支配，由下一节的推导，任意两门课程的相关系数等于二者负荷的乘积：$r_{12}=l_1l_2$，$r_{13}=l_1l_3$，$r_{23}=l_2l_3$。三个方程含三个未知数，可直接求解：
\[
l_1^2=\frac{r_{12}r_{13}}{r_{23}}=\frac{0.83\times0.78}{0.67}=0.966,\qquad l_1=0.983,
\]
\[
l_2=\frac{0.83}{0.983}=0.844,\qquad l_3=\frac{0.78}{0.983}=0.793 .
\]
该结果与\rinline{factanal(covmat=R3,factors=1,n.obs=100)}给出的最大似然负荷（0.983、0.844、0.793）相同，此时模型自由度为0，恰好拟合。变量多于3个时，方程个数多于未知数，“相关系数等于负荷乘积”成为可由数据检验的约束，这是因子模型拟合优度检验的基础。

Spearman正是根据相关系数之间的这种比例关系，推断存在一个一般智力因子。
\end{sikao}
1930年，Louis Thurstone提出多因子理论，并引入因子旋转技术。20世纪50年代至今，随着计算机发展，因子分析应用于心理学以外的众多领域。
\section{因子模型}
\subsection{单因子模型}
学生古典语、法语和英语三门成绩受公因子$f$（语言能力）支配。$f$解释不了的部分称为特殊因子，类似回归中的残差。
\begin{equation}
\begin{cases}
x_1^*=l_1f+u_1,\\
x_2^*=l_2f+u_2,\\
x_3^*=l_3f+u_3.
\end{cases}
\end{equation}
设随机变量$f$和$u_i$满足
\[
\E(f)=0,\quad\var(f)=1,\quad\E(u)=0,\quad
\cov(u)=\Psi=\diag(\psi_1,\psi_2,\psi_3),\quad\cov(u,f)=0.
\]
\Note{$*$表示标准化处理，也可用$x_i-\bar x_i$进行中心化处理。$f$为隐变量或潜变量，需估计$l_1,l_2,l_3$。$\Psi$是对角矩阵，即三个特殊因子互不相关（协方差为0），$\psi_i$表示各自的方差。互不相关只是二阶矩条件；要说相互独立，还需附加条件，例如$(f,u)$服从联合正态分布。}
\begin{figure}[!htbp]\centering
\begin{tikzpicture}[>=Stealth,line width=.65pt,every node/.style={font=\small},every path/.style={draw=Muted}]
\node[draw=Plum,ellipse,text=Plum,minimum width=29mm,minimum height=12mm] (f) at (0,2.5) {语言能力 $f$};
\foreach \i/\lab/\x in {1/古典语/-3,2/法语/0,3/英语/3}{
\node[draw=Ink,minimum width=23mm,minimum height=10mm] (x\i) at (\x,1) {\lab};
\node[draw=Muted,ellipse,fill=PaperTint,minimum width=23mm,minimum height=10mm] (u\i) at (\x,-.5) {残差 $u_\i$};
\draw[->,Plum] (f)--(x\i);\draw[->] (u\i)--(x\i);}
\end{tikzpicture}
\caption{单因子模型}
\end{figure}\FloatBarrier
\subsection{正交多因子模型}
课本例5.8中，Lawley和Maxwell对220名男生的六门课成绩（盖尔语、英语、历史、算术、代数、几何）作因子分析，可能存在两个公因子$f_1$和$f_2$。
\begin{equation}
\begin{aligned}
x_1^*&=l_{11}f_1+l_{12}f_2+u_1,\\
x_2^*&=l_{21}f_1+l_{22}f_2+u_2,\\
&\ \vdots\\
x_6^*&=l_{61}f_1+l_{62}f_2+u_6.
\end{aligned}
\end{equation}
\[
\E(f)=0,\quad\cov(f)=I,\quad\E(u)=0,\quad
\cov(u)=\Psi=\diag(\psi_1,\ldots,\psi_6),\quad\cov(u,f)=0.
\]
\Note{正交两因子模型的$\cov(f)$为$2\times2$单位阵。}
\begin{figure}[!htbp]\centering
\begin{tikzpicture}[>=Stealth,line width=.65pt,every node/.style={font=\small},x=2cm,y=1.15cm]
\node[draw=Plum,text=Plum,ellipse,minimum width=13mm,minimum height=9mm] (f1) at (1.5,2.4) {$f_1$};\node[draw=Forest,text=Forest,ellipse,minimum width=13mm,minimum height=9mm] (f2) at (3.5,2.4) {$f_2$};
\foreach \i in {1,...,6}{\node[draw=Ink,minimum width=13mm,minimum height=8mm] (x\i) at (\i-1,0) {$x_\i^*$};\node (u\i) at (\i-1,-.85) {$u_\i$};\draw[->,Plum!75] (f1)--(x\i);\draw[->,Forest!75,dashed] (f2)--(x\i);\draw[->,Muted] (u\i)--(x\i);}
\end{tikzpicture}
\caption{正交两因子模型}
\end{figure}\FloatBarrier
\subsection{矩阵表示}
以正交两因子模型为例，$X=LF+u$，$L$称为因子负荷系数矩阵（loading matrix），有$p$行$k$列。
\[
X=\begin{bmatrix}x_1^*\\\vdots\\x_6^*\end{bmatrix},\quad
L=\begin{bmatrix}l_{11}&l_{12}\\\vdots&\vdots\\l_{61}&l_{62}\end{bmatrix},\quad
F=\begin{bmatrix}f_1\\f_2\end{bmatrix},\quad
u=\begin{bmatrix}u_1\\\vdots\\u_6\end{bmatrix}.
\]
\[
X=LF+u=
\begin{bmatrix}l_{11}f_1+l_{12}f_2+u_1\\\vdots\\l_{61}f_1+l_{62}f_2+u_6\end{bmatrix}.
\]
\Note{重点是估计出$L$。}
\subsection{负荷矩阵的作用}
求出负荷系数矩阵，就可以写出因子模型表达式。
\begin{equation}
R=R^*+\Psi,\qquad R^*=LL\trans.
\end{equation}
约相关阵$R^*$（Reproduced Correlation Matrix）可还原指标间大部分相关性。

标化变量$X$的协方差阵$\cov(X)$即相关阵$R$。在模型假设下，
\begin{align*}
\cov(X)&=\E\bigl[(LF+u)(LF+u)\trans\bigr]\\
&=\E\bigl[LFF\trans L\trans+uu\trans+LFu\trans+uF\trans L\trans\bigr]\\
&=L\E(FF\trans)L\trans+\E(uu\trans)+L\E(Fu\trans)+\E(uF\trans)L\trans\\
&=LIL\trans+\Psi+0+0=LL\trans+\Psi.
\end{align*}
这是总体模型下的等式。样本中$L$、$\Psi$都要估计，只能近似成立：$R\approx\widehat L\widehat L\trans+\widehat\Psi$。残差阵$E=R-(\widehat L\widehat L\trans+\widehat\Psi)$衡量模型对相关结构的拟合程度。
\Note{$R^*$为$6\times6$矩阵。其第一项对角元素为$l_{11}^2+l_{12}^2$。标化后的方差为1，故$1-(l_{11}^2+l_{12}^2)=\var(u_1)$。例2中，0.634为可被公因子解释的共性方差，0.366由特殊因子解释。}
\section{因子模型的估计方法}
估计方法包括主成分法（principal component，PC）、极大似然法（maximum likelihood，ML）、主因子法（principal factor）、迭代主因子法（iterated principal factor）等。
\Needspace{7\baselineskip}
\subsection{主成分法因子分析}
\Example{例1\quad 三门语言成绩}
从古典语、法语和英语三门成绩的相关矩阵出发，用主成分法作因子分析。
\[
X=\begin{bmatrix}x_1^*\\x_2^*\\x_3^*\end{bmatrix}
=\begin{bmatrix}l_1\\l_2\\l_3\end{bmatrix}f+u,
\qquad
R=LL\trans+\Psi=
\begin{bmatrix}l_1\\l_2\\l_3\end{bmatrix}\begin{bmatrix}l_1&l_2&l_3\end{bmatrix}
+\begin{bmatrix}\psi_1&0&0\\0&\psi_2&0\\0&0&\psi_3\end{bmatrix}.
\]
\[
R=\begin{bmatrix}
l_1^2+\psi_1&l_1l_2&l_1l_3\\
l_1l_2&l_2^2+\psi_2&l_2l_3\\
l_1l_3&l_2l_3&l_3^2+\psi_3
\end{bmatrix}.
\]
\subsubsection{基本思想}
根据矩阵谱分解，相关阵可以用特征值$\lambda_i$和特征向量$e_i$表达：
\[
R=\begin{bmatrix}1&r_{12}&r_{13}\\r_{12}&1&r_{23}\\r_{13}&r_{23}&1\end{bmatrix},\qquad
A=\begin{bmatrix}
\sqrt{\lambda_1}e_{11}&\sqrt{\lambda_2}e_{21}&\sqrt{\lambda_3}e_{31}\\
\sqrt{\lambda_1}e_{12}&\sqrt{\lambda_2}e_{22}&\sqrt{\lambda_3}e_{32}\\
\sqrt{\lambda_1}e_{13}&\sqrt{\lambda_2}e_{23}&\sqrt{\lambda_3}e_{33}
\end{bmatrix}.
\]
\begin{equation}
R=AA\trans=\sum_{i=1}^3\lambda_i e_i e_i\trans.
\end{equation}
视$A$为三因子模型的$L$矩阵，则$R=LL\trans+0$。本例仅一个特征值大于1，且贡献率为84\%，故只保留$A$的第一列，即为主成分解。
\subsubsection{因子分析过程}
\begin{code}
\RFunction{eigen}(R)
\RNumber{2.5219123} \ROperator{/} \RNumber{3}
\end{code}
\[
(\lambda_1,\lambda_2,\lambda_3)=(2.5219123,\ 0.3341080,\ 0.1439797),
\]
\[
\begin{bmatrix}e_1&e_2&e_3\end{bmatrix}=
\begin{bmatrix}
-0.5991079&-0.1148613&0.7923867\\
-0.5729499&-0.6297937&-0.5244885\\
-0.5592836&0.7682231&-0.3115046
\end{bmatrix},\qquad\frac{\lambda_1}{3}=0.8406374.
\]
\[
L=\begin{bmatrix}\sqrt{\lambda_1}e_{11}\\\sqrt{\lambda_1}e_{12}\\\sqrt{\lambda_1}e_{13}\end{bmatrix}
=\begin{bmatrix}1.588\times0.599\\1.588\times0.573\\1.588\times0.559\end{bmatrix}
=\begin{bmatrix}0.951\\0.910\\0.888\end{bmatrix}.
\]
\[
x_1^*=0.951f+u_1,\qquad x_2^*=0.910f+u_2,\qquad x_3^*=0.888f+u_3.
\]
\paragraph{特殊方差与残差矩阵}
主成分法令$\widehat\psi_i=1-\widehat l_i^{2}$，得$\widehat\Psi=\diag(0.095,\ 0.172,\ 0.211)$。残差矩阵为
\[
E=R-(\widehat L\widehat L\trans+\widehat\Psi)=
\begin{bmatrix}0&-0.036&-0.065\\-0.036&0&-0.138\\-0.065&-0.138&0\end{bmatrix}.
\]
以法语与英语为例：观测相关系数为0.67，模型再现$\widehat l_2\widehat l_3=0.910\times0.888=0.808$，残差为$-0.138$。

与Spearman例中的最大似然解（0.983、0.844、0.793）相比，主成分法的负荷（0.951、0.910、0.888）中后两个明显偏大。原因在于主成分法将对角线上的1全部作为待解释的方差，其中本应属于特殊因子的部分也归入了公因子。最大似然解的残差阵为0，而主成分法将法语与英语的相关高估了0.138。

两者反映不同的方面。84\%的贡献率$\lambda_1/3$说明公因子解释了三个标准化变量总方差（对角线）的84\%。残差矩阵看的是非对角的相关系数被再现得怎样。主成分法人为使对角线残差为0，且往往高估负荷，本例三个非对角残差都为负，即一个因子高估了变量间的相关，法语与英语之间尤甚。贡献率高，并不保证相关结构拟合得好。
\subsubsection{SPSS数据格式与语法}
\begin{center}\begin{tabular}{llrrr}\toprule
{\bfseries \texttt{ROWTYPE\_}}&{\bfseries \texttt{VARNAME\_}}&{\bfseries $x_1$}&{\bfseries $x_2$}&{\bfseries $x_3$}\\\midrule
MEAN&&&&\\
STDDEV&&&&\\
N&&100.00&100.00&100.00\\
CORR&$x_1$&1.00&0.83&0.78\\
CORR&$x_2$&0.83&1.00&0.67\\
CORR&$x_3$&0.78&0.67&1.00\\\bottomrule
\end{tabular}\end{center}
\begin{code}
\RKeyword{FACTOR}
 \ROperator{/}\RKeyword{matrix} \RKeyword{in} (cor \ROperator{=} \ROperator{*})
 \ROperator{/}\RKeyword{ANALYSIS} x1 x2 x3
 \ROperator{/}\RKeyword{EXTRACTION} \RKeyword{PC}
 \ROperator{/}\RKeyword{ROTATION} \RKeyword{NOROTATE}
\end{code}
\Example{例2\quad 220名男生六门课成绩（课本例5.8）}
要求：确定因子个数，并用主成分法作因子分析。
\par\addvspace{5pt}\begin{center}\begin{minipage}{\linewidth}\centering
\small
\captionof{table}{六门课成绩的相关矩阵}\label{tab:six}
\begin{tabular}{lrrrrrr}\toprule
 &{\bfseries  盖尔语 }&{\bfseries  英语 }&{\bfseries  历史 }&{\bfseries  算术 }&{\bfseries  代数 }&{\bfseries  几何}\\\midrule
盖尔语 & 1.000 & 0.439 & 0.410 & 0.288 & 0.329 & 0.248\\
英语 & 0.439 & 1.000 & 0.351 & 0.354 & 0.320 & 0.329\\
历史 & 0.410 & 0.351 & 1.000 & 0.164 & 0.190 & 0.181\\
算术 & 0.288 & 0.354 & 0.164 & 1.000 & 0.595 & 0.470\\
代数 & 0.329 & 0.320 & 0.190 & 0.595 & 1.000 & 0.464\\
几何 & 0.248 & 0.329 & 0.181 & 0.470 & 0.464 & 1.000\\
\bottomrule\end{tabular}
\end{minipage}\end{center}

\subsubsection{因子负荷系数矩阵的性质}
\begin{enumerate}
\item 行元素平方和即共性方差，说明变量主要受公因子还是特殊因子支配。
\item 列元素平方和即因子贡献。
\item 观测变量标准化（$\var X_i=1$）、因子方差为1且因子正交（$\cov(F)=I$）、$\cov(F,u)=0$时，负荷系数即因子与变量间的相关系数：
\[
\cov(X_i,F_j)=\cov\Bigl(\sum_{s=1}^kl_{is}F_s+u_i,\ F_j\Bigr)=l_{ij},\qquad
\operatorname{corr}(X_i,F_j)=\frac{l_{ij}}{\sqrt{\var X_i\var F_j}}=l_{ij}.
\]
变量未标准化时，相关系数为$l_{ij}/\sigma_i$；因子斜交时，$\cov(X_i,F_j)$为$L\Phi$的元素（$\Phi$为因子相关阵），不再等于$l_{ij}$。
\end{enumerate}
本例相关阵的特征值为2.733、1.130、0.615、0.601、0.525、0.396，有两个大于1，按“特征值大于1”保留2个因子。这只是经验准则（Kaiser准则），还应结合碎石图、拟合检验、残差和因子能否解释来定。
\subsection{最大似然法因子分析}
\Example{例3\quad 血压与体型指标}
测定742名正常人的收缩压（mmHg）、舒张压（mmHg）、体重（kg）、腰围（cm）和臀围（cm）五个指标。假设它们服从均数为$\mu$、协方差阵为$\Sigma$的多元正态分布，要求从相关阵出发用最大似然法作因子分析。原始数据为课程文件\texttt{exp2forMLE.sav}。
\[
X=LF+u,\qquad \Sigma=\cov(X)=LL\trans+\Psi,\qquad
\Psi=\diag(\psi_1,\ldots,\psi_5).
\]
\[
X\sim N_p(\mu,\Sigma),\qquad\mathcal L=\prod_{i=1}^n f(x_i),\qquad\log\mathcal L\ \text{含未知参数}.
\]
\Note{要求原变量符合多元正态分布。}
\subsubsection{正态分布}
设$x$服从均数为$\mu$、方差为$\sigma^2$的一元正态分布，则
\[
f(x)=\frac{1}{(2\pi)^{1/2}\sigma}\exp\left[-\frac12(x-\mu)(\sigma^2)^{-1}(x-\mu)\right]=\frac{1}{\sqrt{2\pi}\,\sigma}\exp\left[-\frac{(x-\mu)^2}{2\sigma^2}\right].
\]
设随机向量$X=(x_1,\ldots,x_p)\trans$服从$p$元正态分布，则
\begin{equation}
f(X)=\frac{1}{(2\pi)^{p/2}|\Sigma|^{1/2}}
\exp\left[-\frac12(X-\mu)\trans\Sigma^{-1}(X-\mu)\right].
\end{equation}
$\mu$为均数向量，$\Sigma$为$x_1,\ldots,x_p$的协方差阵，$|\Sigma|$为其行列式值。
\subsubsection{一元最大似然估计}
\begin{code}
data \ROperator{<-} \RFunction{rnorm}(\RNumber{200}, \RNumber{10}, \RNumber{2})
nll1 \ROperator{<-} \RKeyword{function}(params) \{
  mu \ROperator{<-} params[\RNumber{1}]
  sigma \ROperator{<-} \RFunction{exp}(params[\RNumber{2}])          \RComment{\# 对数尺度参数化，保证sigma > 0}
  \ROperator{-}\RFunction{sum}(\RFunction{dnorm}(data, mean \ROperator{=} mu, sd \ROperator{=} sigma, log \ROperator{=} \RKeyword{TRUE}))
\}
fit \ROperator{<-} \RFunction{optim}(par \ROperator{=} \RFunction{c}(\RNumber{8}, \RFunction{log}(\RNumber{3})), fn \ROperator{=} nll1, method \ROperator{=} \RString{"BFGS"})
\RFunction{c}(mu \ROperator{=} fit\ROperator{\$}par[\RNumber{1}], sigma \ROperator{=} \RFunction{exp}(fit\ROperator{\$}par[\RNumber{2}]))
\end{code}
估计值与闭式解$\bar x$及$\sqrt{\sum(x_i-\bar x)^2/n}$一致（注意分母为$n$）。
\subsubsection{多元最大似然估计}
\begin{code}
\RFunction{library}(MASS)                      \RComment{\# mvrnorm() 在 MASS 包中}
\RFunction{library}(mvtnorm)                   \RComment{\# dmvnorm() 在 mvtnorm 包中}
mu \ROperator{<-} \RFunction{c}(\RNumber{152}, \RNumber{38})
sigma \ROperator{<-} \RFunction{matrix}(\RFunction{c}(\RNumber{70}, \RNumber{34}, \RNumber{34}, \RNumber{30}), \RNumber{2}, \RNumber{2})
da \ROperator{<-} \RFunction{mvrnorm}(\RNumber{100}, mu, sigma)
nll2 \ROperator{<-} \RKeyword{function}(params) \{
  m \ROperator{<-} params[\RNumber{1}\ROperator{:}\RNumber{2}]                       \RComment{\# 均值向量}
  U \ROperator{<-} \RFunction{matrix}(\RFunction{c}(\RFunction{exp}(params[\RNumber{3}]), \RNumber{0}, params[\RNumber{4}], \RFunction{exp}(params[\RNumber{5}])), \RNumber{2}, \RNumber{2})
  S \ROperator{<-} \RFunction{crossprod}(U)                     \RComment{\# S = t(U) U，必为正定}
  \ROperator{-}\RFunction{sum}(\RFunction{dmvnorm}(da, mean \ROperator{=} m, sigma \ROperator{=} S, log \ROperator{=} \RKeyword{TRUE}))
\}
U0 \ROperator{<-} \RFunction{chol}(\RFunction{cov}(da))                      \RComment{\# 以样本矩估计作初值}
par0 \ROperator{<-} \RFunction{c}(\RFunction{colMeans}(da), \RFunction{log}(U0[\RNumber{1}, \RNumber{1}]), U0[\RNumber{1}, \RNumber{2}], \RFunction{log}(U0[\RNumber{2}, \RNumber{2}]))
mle\_result \ROperator{<-} \RFunction{optim}(par \ROperator{=} par0, fn \ROperator{=} nll2, method \ROperator{=} \RString{"BFGS"})
\end{code}
二元正态有2个均值和3个协方差参数（$\sigma_{11},\sigma_{12},\sigma_{22}$），共5个未知参数。上面用Cholesky分解$\Sigma=U\trans U$参数化（$U$为上三角阵，对角元取指数），恰好5个自由参数，且对任何参数值$\Sigma$都正定，优化不会走到不合法的区域。

设$x_1,\ldots,x_n$独立同分布于$N_p(\mu,\Sigma)$。不论$\Sigma$取何值，$\mu$的极大似然估计都是$\bar x$；代入后得到只含$\Sigma$的对数似然。记$S_n=\frac1n\sum_{i=1}^n(x_i-\bar x)(x_i-\bar x)\trans$（分母为$n$），则
\begin{align*}
\log\mathcal L(L,\Psi)
&=C-\frac n2\log|\Sigma|-\frac12\tr(\Sigma^{-1}nS_n)\\
&=C-\frac n2\left\{\log|LL\trans+\Psi|+\tr\bigl[(LL\trans+\Psi)^{-1}S_n\bigr]\right\}.
\end{align*}
不加约束时，最大值在$\Sigma=S_n$处取得。因子分析是在$\Sigma=LL\trans+\Psi$（$L$为$p\times k$阵，$\Psi$为正对角阵）的约束下求极大似然，没有闭式解，只能迭代。ML估计具有尺度不变性，从相关阵或协方差阵出发得到等价的解。
\subsubsection{模型自由度}
以$p$为变量数、$k$为因子数。$\Sigma$（或$R$）有$p(p+1)/2$个不同元素。模型参数为$L$的$pk$个元素和$\Psi$的$p$个对角元；因子方差已固定为1、因子间协方差固定为0（$\cov(F)=I$），不再计为参数。但对任意$k\times k$正交阵$T$，$L^*=LT$与$L$给出同一个$LL\trans$（见\ref{sec:rotation}节），模型无法区分；$k\times k$正交阵有$k(k-1)/2$个自由度，必须扣除（估计时常加“$L\trans\Psi^{-1}L$为对角阵”之类的约束来取定唯一解）。于是
\begin{equation}
df=\frac{p(p+1)}2-\left[pk+p-\frac{k(k-1)}2\right]=\frac{(p-k)^2-p-k}{2}.\label{eq:fa-df}
\end{equation}
展开即可验证两式相等：$\bigl[p^2+p-2pk-2p+k^2-k\bigr]/2=\bigl[(p-k)^2-p-k\bigr]/2$。六门课例中$p=6$、$k=2$：$df=21-(12+6-1)=4$。$df>0$时模型对$\Sigma$施加了实质约束，才能检验拟合优度。
\Note{有的资料把参数数记为$p(k+1)+k$，即把因子方差也算作参数，且没有扣除旋转不确定性。按此计算六门课例得$df=21-20=1$，与上式的4矛盾，应以上式为准。}
\subsubsection{拟合优度检验}
最大似然法因子分析，当模型自由度大于0时，可进行拟合优度检验。
\[
H_0:\Sigma=LL\trans+\Psi,\qquad H_1:\Sigma\text{为其他正定矩阵},\qquad \alpha=0.05.
\]
\begin{equation}
\chi^2=\left(n-1-\frac{2p+4k+5}{6}\right)\log\frac{|\widehat\Sigma|}{|S_n|}\ \dot\sim\ \chi^2(df),
\end{equation}
$df$按式\eqref{eq:fa-df}计算。从相关阵出发调用\texttt{factanal()}时，必须用\texttt{n.obs}给出样本量，否则无法计算检验统计量：
\begin{code}
r \ROperator{<-} \RFunction{c}(\RNumber{.439, .410, .288, .329, .248, .351, .354, .320, .329, .164, .190, .181, .595, .470, .464})
R \ROperator{<-} \RFunction{diag}(\RNumber{6}); R[\RFunction{lower.tri}(R)] \ROperator{<-} r; R \ROperator{<-} R \ROperator{+} \RFunction{t}(R) \ROperator{-} \RFunction{diag}(\RNumber{6})   \RComment{\# 表\ref{tab:six}}
out1 \ROperator{<-} \RFunction{factanal}(covmat \ROperator{=} R, factors \ROperator{=} \RNumber{2}, n.obs \ROperator{=} \RNumber{220}, rotation \ROperator{=} \RString{"none"})
out1\ROperator{\$}STATISTIC; out1\ROperator{\$}dof; out1\ROperator{\$}PVAL
\end{code}
完整输出如下（直接采用varimax旋转，旋转不影响检验结果）：
\par\noindent\begin{rcode}
out2 <- factanal(covmat=R,factors=2,n.obs=220,rotation="varimax")
print(out2, cutoff=0)
\end{rcode}
\begin{routput}
Uniquenesses:
  Gaelic  English  History    Arith  Algebra Geometry 
   0.510    0.594    0.644    0.377    0.431    0.628 

Loadings:
         Factor1 Factor2
Gaelic   0.235   0.659  
English  0.323   0.549  
History  0.088   0.590  
Arith    0.771   0.170  
Algebra  0.724   0.213  
Geometry 0.572   0.210  

               Factor1 Factor2
SS loadings      1.612   1.203
Proportion Var   0.269   0.201
Cumulative Var   0.269   0.469

Test of the hypothesis that 2 factors are sufficient.
The chi square statistic is 2.33 on 4 degrees of freedom.
The p-value is 0.674 
\end{routput}
输出各项的含义：\rinline{Uniquenesses}为特殊方差$\widehat\psi_i$，$1-\widehat\psi_i$为共性方差（如算术$1-0.377=0.623$）；\rinline{Loadings}为负荷矩阵；\rinline{SS loadings}为各列平方和，即因子贡献；\rinline{Proportion Var}为因子贡献除以变量数6。最后两行为拟合优度检验。

仅取一个因子时，同一检验给出$\chi^2=52.0$，$df=9$，$P=4.5\times10^{-8}$，拒绝单因子模型。单一的一般能力因子不足以解释六门课程之间的相关，需区分数学与文史两个方面。此例说明拟合优度检验可辅助确定因子个数。

六门课例：$\chi^2=2.33$，$df=4$，$P=0.674$。$P>0.05$表示在$\alpha=0.05$水平上未拒绝两因子的协方差结构，并不等于证明模型正确。样本量大时微小偏离也会被拒绝，样本量小时检验效能又低。本例残差相关的最大绝对值为0.030，各变量共性方差在0.36～0.62之间，两个因子（数学、文史）都能解释，几方面结合起来看，两因子模型是合适的。
\subsection{主因子法与迭代主因子法}
对例3的五项指标，从相关阵出发用主因子法作因子分析。
\begin{enumerate}
\item 估计共性方差初值。常用复相关系数平方$h_i^{2(0)}=1-1/r^{ii}$（$r^{ii}$为$R^{-1}$的第$i$个对角元，即$x_i$对其余变量回归的$R^2$），也可用该变量与其他变量相关系数绝对值的最大值。
\item 用共性方差初值代替相关阵的主对角线，得到约相关阵$R_h=R-\widehat\Psi$。
\item 对约相关阵求特征值和特征向量，取前$k$个，$\widehat L=(\sqrt{\lambda_1}e_1,\ldots,\sqrt{\lambda_k}e_k)$，即主因子解。
\item 迭代：令$h_i^2=\sum_{j=1}^k\widehat l_{ij}^{2}$、$\widehat\psi_i=1-h_i^2$，回到第2步，直至相邻两次的$\max_i|h_i^{2(t+1)}-h_i^{2(t)}|$小于容许误差（如$10^{-4}$）或达到最大迭代次数。
\end{enumerate}
\Note{迭代主因子法收敛时，等价于使非对角残差平方和最小的（未加权）最小二乘解。最大似然法则在正态似然下按$1/\psi_i$加权。两者相似但不相同，都可能出现下述Heywood现象。}
\subsubsection{Heywood现象}
用最大似然法或主因子法作因子分析时，有时会出现共性方差$h_i^2\geqslant1$的情况，称为Heywood现象，给解释带来困难。应区分两种情形。
\begin{itemize}
\item $\widehat\psi_i=0$（$h_i^2=1$）：边界解，估计落在参数空间的边界上，该变量被公因子完全解释，相应的标准误和检验不可靠。R的\texttt{factanal()}把$\psi_i$限制为不小于0.005，遇到这种情形会停在下界。
\item $\widehat\psi_i<0$（$h_i^2>1$）：方差为负，已在参数空间之外，是不容许解（improper solution），不能解释。
\end{itemize}
常见原因有样本量小、因子数取得过多或过少、模型设定不当、个别变量与其他变量高度共线等。
\Note{课本72页。}
\section{因子旋转}\label{sec:rotation}
上述方法求出的是因子分析的初始解。有时初始因子不易解释，实际应用中常对负荷作正交变换，即旋转，以得到更简单的结构，便于因子解释。最常用的方法是方差最大（varimax）正交旋转。

在保持因子间正交、各变量共性方差不变的条件下，使负荷结构尽量简单：每一列的平方负荷尽量两极分化，接近0或接近1。采用Kaiser标准化时，令$h_i^2=\sum_{j=1}^kl_{ij}^2$，规范化负荷为$\tilde l_{ij}=l_{ij}/h_i$，即把每行化为单位长度。varimax选择正交阵$T$，使各列规范化平方负荷的方差之和
\[
V=\sum_{j=1}^k\left[\frac1p\sum_{i=1}^p\tilde l_{ij}^{\,4}-\Bigl(\frac1p\sum_{i=1}^p\tilde l_{ij}^{\,2}\Bigr)^2\right]
\]
达到最大，求得$T$后再把各行乘回$h_i$。R的\texttt{varimax()}默认\texttt{normalize = TRUE}，即采用Kaiser标准化。从几何观点看，正交旋转保持原坐标轴的正交性，对其进行旋转。
\Example{例5\quad 六门课成绩的方差最大正交旋转}
利用220名男生六门课成绩的相关矩阵，用最大似然法作因子分析，并作方差最大正交旋转。
\begin{code}
out1 \ROperator{<-} \RFunction{factanal}(covmat \ROperator{=} R, factors \ROperator{=} \RNumber{2}, n.obs \ROperator{=} \RNumber{220}, rotation \ROperator{=} \RString{"none"})
\RFunction{loadings}(out1)
out2 \ROperator{<-} \RFunction{factanal}(covmat \ROperator{=} R, factors \ROperator{=} \RNumber{2}, n.obs \ROperator{=} \RNumber{220}, rotation \ROperator{=} \RString{"varimax"})
\RFunction{print}(\RFunction{loadings}(out2), cutoff \ROperator{=} \RNumber{0})   \RComment{\# 默认 cutoff = 0.1 会隐藏小负荷}
\end{code}
\par\addvspace{5pt}\begin{center}\begin{minipage}{\linewidth}\centering
\small
\captionof{table}{旋转前后的负荷矩阵及方差贡献}\label{tab:rot}
\begin{tabular}{lrrrr}\toprule
{\bfseries 指标 }&{\bfseries  初始 $f_1$ }&{\bfseries  初始 $f_2$ }&{\bfseries  旋转 $f_1$ }&{\bfseries  旋转 $f_2$}\\\midrule
盖尔语 & 0.553 & 0.429 & 0.235 & 0.659\\
英语 & 0.568 & 0.288 & 0.323 & 0.549\\
历史 & 0.392 & 0.450 & 0.088 & 0.590\\
算术 & 0.740 & -0.273 & 0.771 & 0.170\\
代数 & 0.724 & -0.211 & 0.724 & 0.213\\
几何 & 0.595 & -0.132 & 0.572 & 0.210\\
载荷平方和 & 2.209 & 0.606 & 1.612 & 1.203\\
方差贡献比例 & 0.368 & 0.101 & 0.269 & 0.201\\
累积贡献比例 & 0.368 & 0.469 & 0.269 & 0.469\\
\bottomrule\end{tabular}
\end{minipage}\end{center}

\Note{旋转后，$f_1$为数学因子，$f_2$为文史因子。表中历史在旋转后$f_1$上的负荷0.088，在默认输出中因\texttt{print()}的\texttt{cutoff = 0.1}而显示为空白，此处补出。}
\begin{zm}[旋转不改变拟合]
设$T$为正交阵，$L^*=LT$，则
\[
L^*L^{*\mathsf T}=LTT\trans L\trans=LL\trans.
\]
所以$\widehat\Sigma=LL\trans+\Psi$、残差阵、似然值和$\chi^2$检验都不变；$LL\trans$的对角元即共性方差$h_i^2$，也不变。
\end{zm}
各列平方和（因子贡献）则可以改变，即公因子解释的方差在因子间重新分配。本例总共性方差$2.209+0.606=1.612+1.203=2.815$保持不变，旋转后两个因子的贡献更均衡。

\begin{sikao}[因子旋转的依据]
同一资料经旋转后负荷矩阵发生变化，旋转前后的两组负荷均为同一模型的有效表示。

因子不可观测，模型只规定了$LL\trans$，即变量间的相关结构，而未规定坐标轴的方向。这与描述平面上点的位置类似：采用东西--南北坐标或旋转$30^\circ$后的坐标均可，点的位置不变。本例旋转前，第一因子是与所有课程相关的一般因子，第二因子是文史与数学的对比；旋转后，两个因子分别对应数学课程与文史课程。两种表示的似然值、$\chi^2$与残差完全相同，选择的依据是何者更便于解释。

因子的名称由分析者依据负荷确定，具有一定的主观性。因子分析的结论宜通过新资料或专业理论加以验证。
\end{sikao}
求varimax的正交阵$T$一般只能用迭代法。当$k=2$时，可以精确求出$2\times2$正交变换矩阵$T$：
\[
T_{\text{顺时针}}=\begin{bmatrix}\cos\theta&\sin\theta\\-\sin\theta&\cos\theta\end{bmatrix},\qquad
T_{\text{逆时针}}=\begin{bmatrix}\cos\theta&-\sin\theta\\\sin\theta&\cos\theta\end{bmatrix}.
\]
\begin{code}
out2\ROperator{\$}rotmat
\RFunction{acos}(\RNumber{0.8420397}) \ROperator{/} pi \ROperator{*} \RNumber{180}
\RFunction{asin}(\RNumber{0.5394156}) \ROperator{/} pi \ROperator{*} \RNumber{180}
\end{code}
\[
T=\begin{bmatrix}0.8420397&0.5394156\\-0.5394156&0.8420397\end{bmatrix},\qquad
\theta=32.64386^\circ\quad(\arccos),\quad32.64387^\circ\quad(\arcsin).
\]
\begin{figure}[!htbp]\centering
\begin{tikzpicture}\begin{axis}[width=9.6cm,height=8cm,xmin=-1.05,xmax=1.05,ymin=-1.05,ymax=1.05,axis equal image,xlabel={$L_{\cdot1}$},ylabel={$L_{\cdot2}$},xtick={-1,-.5,0,.5,1},ytick={-1,-.5,0,.5,1}]\addplot[only marks,mark=*,mark size=2.6pt,Ink] coordinates {(0.553,0.429) (0.568,0.288) (0.392,0.45) (0.74,-0.273) (0.724,-0.211) (0.595,-0.132)};\addplot[Muted,line width=.6pt,->] coordinates {(-1,0) (1,0)};\addplot[Muted,line width=.6pt,->] coordinates {(0,-1) (0,1)};\addplot[Brick,densely dashed,line width=1pt] coordinates {(-1,.64060594) (1,-.64060594)};\addplot[Forest,densely dashed,line width=1pt] coordinates {(-.64060594,-1) (.64060594,1)};\end{axis}\end{tikzpicture}
\caption{六门课因子负荷与正交旋转（$32.64386^\circ$）}
\end{figure}\FloatBarrier

\section{因子得分}
公共因子的估计，即因子得分，是对不可观测、抽象但有实际意义的随机变量的估计。
\subsection{估计方法}
\paragraph{Bartlett法} 给定$L$和$\Psi$，把个体的因子值$f$当作待估参数，使加权残差平方和
\[
Q(f)=(x^*-Lf)\trans\Psi^{-1}(x^*-Lf)
\]
最小，权重为特殊方差的倒数。一阶条件为$\partial Q/\partial f=-2L\trans\Psi^{-1}(x^*-Lf)=0$，解得
\begin{equation}
\widehat f_b=(L\trans\Psi^{-1}L)^{-1}L\trans\Psi^{-1}x^*.
\end{equation}
这就是以$\Psi^{-1}$为权的加权最小二乘回归。若进一步假定$x^*\mid f\sim N(Lf,\Psi)$，$Q(f)$就是$-2\times$对数似然（相差常数），$\widehat f_b$也是$f$的极大似然估计。这一“极大似然”解释仅限于给定$L$、$\Psi$的正态条件模型。它条件无偏：$\E(\widehat f_b\mid f)=f$。
\paragraph{Thomson法（回归法）} 把$f$当作随机变量，$\E f=0$，$\cov(f)=I$，$\cov(f,x^*)=L\trans$，$\cov(x^*)=\Sigma=LL\trans+\Psi$。使$\E\|f-Ax^*\|^2$最小的线性预测为$A=\cov(f,x^*)\cov(x^*)^{-1}=L\trans\Sigma^{-1}$，即
\begin{equation}
\widehat f_t=L\trans\Sigma^{-1}x^*=(I+L\trans\Psi^{-1}L)^{-1}L\trans\Psi^{-1}x^*.
\end{equation}
\begin{zm}[两种写法相等]
$(I+L\trans\Psi^{-1}L)L\trans=L\trans+L\trans\Psi^{-1}LL\trans=L\trans\Psi^{-1}(\Psi+LL\trans)=L\trans\Psi^{-1}\Sigma$。左乘$(I+L\trans\Psi^{-1}L)^{-1}$、右乘$\Sigma^{-1}$即得。
\end{zm}
$\widehat f_t$是最优线性预测；若$(f,x^*)$服从联合正态分布，它就是条件期望$\E(f\mid x^*)$，预测误差的协方差阵为$\cov(f-\widehat f_t)=I-L\trans\Sigma^{-1}L$。Thomson得分向0收缩，条件有偏，但均方误差比Bartlett得分小；Bartlett得分条件无偏，方差较大。
\Note{可以使用旋转前或旋转后的因子。无论哪种方法，因子得分都只是预测值，含有预测误差，并不是个体真实的因子值；实际计算时$L$、$\Psi$又用估计值代替。把得分当作无误差的变量再放进回归等分析，会低估不确定性。}
计算个体得分需要每个人的标准化观测值$x^*$。只有相关矩阵时，只能算出得分系数$L\trans R^{-1}$，算不出个体得分；R中\texttt{factanal(x = 原始数据, scores = "regression")}可以给出得分，\texttt{factanal(covmat = R)}则不能。
\Example{例6\quad 计算因子得分}
测定742名正常人的收缩压、舒张压、体重、腰围和臀围五个指标，从相关阵出发用主成分法作因子分析，作方差最大正交旋转，并计算回归法因子得分。原始数据为课程文件\texttt{exp2forMLE.sav}（742行，变量\texttt{ssy}、\texttt{szy}、\texttt{tz}、\texttt{yw}、\texttt{tw}），例3也用这份数据。
\begin{code}
da \ROperator{<-} foreign\ROperator{::}\RFunction{read.spss}(\RString{"exp2forMLE.sav"}, to.data.frame \ROperator{=} \RKeyword{TRUE})
R \ROperator{<-} \RFunction{cor}(da); e \ROperator{<-} \RFunction{eigen}(R)                \RComment{\# 特征值 3.09, 1.25, 0.34, ...}
L \ROperator{<-} e\ROperator{\$}vectors[, \RNumber{1}\ROperator{:}\RNumber{2}] \ROperator{\%*\%} \RFunction{diag}(\RFunction{sqrt}(e\ROperator{\$}values[\RNumber{1}\ROperator{:}\RNumber{2}]))   \RComment{\# 主成分解}
L \ROperator{<-} \RFunction{unclass}(\RFunction{varimax}(L)\ROperator{\$}loadings)          \RComment{\# 方差最大旋转}
F \ROperator{<-} \RFunction{scale}(da) \ROperator{\%*\%} \RFunction{solve}(R) \ROperator{\%*\%} L         \RComment{\# 回归法：每行为 t(L) R^\{-1\} x*}
\end{code}
旋转后第一因子在体重、腰围、臀围上的负荷为0.93～0.95，为体型因子；第二因子在收缩压、舒张压上的负荷约0.88～0.89，为血压因子。

\begin{chaptersummary}{本章小结：因子分析}
研究目的&用少数不可观测的公因子解释多个观测变量之间的相关（协方差结构），并给公因子以专业解释\\
统计模型&$x^*=Lf+u$，$\cov(f)=I$，$\cov(u)=\Psi$为对角阵，$\cov(f,u)=0$；于是$\Sigma=LL\trans+\Psi$\\
主要条件&变量间有相关且存在共同因子结构（KMO、Bartlett检验作参考）；ML法及其拟合检验要求多元正态、样本量足够；$df>0$模型才可检验\\
结果解释&负荷（标准化、正交时即相关系数）、共性方差$h_i^2$、因子贡献、残差相关阵；旋转只重新分配贡献，不改变拟合\\
常见误用&把特殊因子“互不相关”说成“独立”；只凭特征值大于1或$P>0.05$定因子数；忽视Heywood解；把84\%之类的贡献率当作相关结构拟合良好；以为只有相关矩阵也能算个体得分；把得分当作无误差的真值\\
\end{chaptersummary}

\begin{fuxi}[复习题]
\begin{enumerate}
\item 标准化变量的正交因子模型中，某变量在两个因子上的负荷为0.6和0.5，计算其共性方差与特殊方差。\\
\textit{答：共性方差$0.36+0.25=0.61$，特殊方差$0.39$。}
\item $p=6$个变量取$k=2$个因子，计算最大似然因子分析拟合优度检验的自由度。\\
\textit{答：$[(6-2)^2-6-2]/2=4$。}
\item 列出因子旋转前后保持不变的量与发生变化的量。\\
\textit{答：共性方差、$LL\trans$、残差阵、似然值与$\chi^2$检验不变；负荷矩阵和各因子贡献（列平方和）会变，总贡献不变。}
\item 主成分法因子分析中第一因子的贡献率为84\%，说明这一数值能否用于判断单因子模型的拟合程度。\\
\textit{答：不能。贡献率反映对角线上的方差，拟合程度需看非对角线的残差相关是否足够小；采用最大似然法时还可作$\chi^2$检验。}
\item 比较主成分分析与因子分析的根本区别。\\
\textit{答：主成分分析是变量的线性变换，不设模型；因子分析是统计模型，假定公因子导致变量相关，并单独设特殊因子（误差）。}
\end{enumerate}
\end{fuxi}
